\documentclass[10pt,a4paper]{amsart}
\usepackage{amsmath,amssymb,amsthm,graphicx}
% Retained preamble for the Stacks Project chapter snapshot in this directory.
% The chapter's own preamble.tex is a copy of the shared Stacks Project preamble; the block
% below is sliced MECHANICALLY, line for line, from it. Nothing was added, removed or
% redefined; the wrapper class, the xr-hyper/hyperref wiring and the externaldocument list
% are omitted because the assembler supplies the document class.
%
% lines taken from preamble.tex: 3, 5, 9, 11, 12, 156, 157, 158, ... (51 lines total)
%
% package coverage restored from the chapter's own preamble (2026-10-13): xy, xy-2cell

\usepackage[all]{xy}
\xyoption{2cell}
\UseAllTwocells

\newenvironment{reference}{\comment}{\endcomment}
\newenvironment{slogan}{\comment}{\endcomment}
\newenvironment{history}{\comment}{\endcomment}
\theoremstyle{plain}
\newtheorem{theorem}[subsection]{Theorem}
\newtheorem{proposition}[subsection]{Proposition}
\newtheorem{lemma}[subsection]{Lemma}
\theoremstyle{definition}
\newtheorem{definition}[subsection]{Definition}
\newtheorem{example}[subsection]{Example}
\newtheorem{exercise}[subsection]{Exercise}
\newtheorem{situation}[subsection]{Situation}
\theoremstyle{remark}
\newtheorem{remark}[subsection]{Remark}
\newtheorem{remarks}[subsection]{Remarks}
\numberwithin{equation}{subsection}
\def\lim{\mathop{\mathrm{lim}}\nolimits}
\def\colim{\mathop{\mathrm{colim}}\nolimits}
\def\Spec{\mathop{\mathrm{Spec}}}
\def\Hom{\mathop{\mathrm{Hom}}\nolimits}
\def\Ext{\mathop{\mathrm{Ext}}\nolimits}
\def\SheafHom{\mathop{\mathcal{H}\!\mathit{om}}\nolimits}
\def\SheafExt{\mathop{\mathcal{E}\!\mathit{xt}}\nolimits}
\def\Sch{\mathit{Sch}}
\def\Mor{\mathop{\mathrm{Mor}}\nolimits}
\def\Ob{\mathop{\mathrm{Ob}}\nolimits}
\def\Sh{\mathop{\mathit{Sh}}\nolimits}
\def\NL{\mathop{N\!L}\nolimits}
\def\CH{\mathop{\mathrm{CH}}\nolimits}
\def\proetale{{pro\text{-}\acute{e}tale}}
\def\etale{{\acute{e}tale}}
\def\QCoh{\mathit{QCoh}}
\def\Ker{\mathop{\mathrm{Ker}}}
\def\Im{\mathop{\mathrm{Im}}}
\def\Coker{\mathop{\mathrm{Coker}}}
\def\Coim{\mathop{\mathrm{Coim}}}
\DeclareMathSymbol{\boxtimes}{\mathbin}{AMSa}{"02}
\def\QCohstack{\mathcal{QC}\!\mathit{oh}}
\def\Cohstack{\mathcal{C}\!\mathit{oh}}
\def\Spacesstack{\mathcal{S}\!\mathit{paces}}
\def\Quotfunctor{\mathrm{Quot}}
\def\Hilbfunctor{\mathrm{Hilb}}
\def\Curvesstack{\mathcal{C}\!\mathit{urves}}
\def\Polarizedstack{\mathcal{P}\!\mathit{olarized}}
\def\Complexesstack{\mathcal{C}\!\mathit{omplexes}}
\def\Pic{\mathop{\mathrm{Pic}}\nolimits}
\def\Picardstack{\mathcal{P}\!\mathit{ic}}
\def\Picardfunctor{\mathrm{Pic}}
\def\Deformationcategory{\mathcal{D}\!\mathit{ef}}
% lane-s11: discard bibliographic \begin{reference} blocks entirely (incl. any \cite inside).
\newsavebox{\limrefbox}
\ifcsname reference\endcsname\expandafter\let\csname reference\endcsname\undefined\fi\ifcsname endreference\endcsname\expandafter\let\csname endreference\endcsname\undefined\fi
\def\citation#1{\relax}
\newenvironment{reference}{\setbox\limrefbox=\vbox\bgroup\let\cite\relax}{\egroup\ignorespaces}
\usepackage{cleveref}
\crefname{originalproposition}{original proposition}{original propositions}
\Crefname{originalproposition}{Original proposition}{Original propositions}
\makeatletter
\newlabel{proposition-approximate}{{01ZA}{source}{}{sourcepointer-1}{}}
\newlabel{proposition-approximate@cref}{{[originalproposition][0][]01ZA}{[1][0][]source}{}{}{}}
\makeatother
\crefname{originalreference}{original reference}{original references}
\Crefname{originalreference}{Original reference}{Original references}
\makeatletter
\newlabel{morphisms-lemma-noetherian-finite-type-finite-presentation}{{01TX}{source}{}{sourcepointer-2}{}}
\newlabel{morphisms-lemma-noetherian-finite-type-finite-presentation@cref}{{[originalreference][0][]01TX}{[1][0][]source}{}{}{}}
\makeatother
\crefname{originalreference}{original reference}{original references}
\Crefname{originalreference}{Original reference}{Original references}
\makeatletter
\newlabel{morphisms-lemma-base-change-finite-presentation}{{01TS}{source}{}{sourcepointer-3}{}}
\newlabel{morphisms-lemma-base-change-finite-presentation@cref}{{[originalreference][0][]01TS}{[1][0][]source}{}{}{}}
\makeatother
\crefname{originalreference}{original reference}{original references}
\Crefname{originalreference}{Original reference}{Original references}
\makeatletter
\newlabel{morphisms-lemma-locally-finite-type-characterize}{{01T2}{source}{}{sourcepointer-4}{}}
\newlabel{morphisms-lemma-locally-finite-type-characterize@cref}{{[originalreference][0][]01T2}{[1][0][]source}{}{}{}}
\makeatother
\crefname{originallemma}{original lemma}{original lemmas}
\Crefname{originallemma}{Original lemma}{Original lemmas}
\makeatletter
\newlabel{lemma-descend-opens}{{01Z4}{source}{}{sourcepointer-5}{}}
\newlabel{lemma-descend-opens@cref}{{[originallemma][0][]01Z4}{[1][0][]source}{}{}{}}
\makeatother
\crefname{originallemma}{original lemma}{original lemmas}
\Crefname{originallemma}{Original lemma}{Original lemmas}
\makeatletter
\newlabel{lemma-limit-affine}{{01Z6}{source}{}{sourcepointer-6}{}}
\newlabel{lemma-limit-affine@cref}{{[originallemma][0][]01Z6}{[1][0][]source}{}{}{}}
\makeatother
\crefname{originallemma}{original lemma}{original lemmas}
\Crefname{originallemma}{Original lemma}{Original lemmas}
\makeatletter
\newlabel{lemma-descend-section}{{01Z0}{source}{}{sourcepointer-7}{}}
\newlabel{lemma-descend-section@cref}{{[originallemma][0][]01Z0}{[1][0][]source}{}{}{}}
\makeatother
\title{Locally finite type morphisms occur inside finite presentation approximations}
\author{The Stacks Project Authors}
\date{stacks-project:limits.tex:local-snapshot-20261012; GNU FDL 1.2 or later}
\begin{document}
\maketitle
\noindent\textit{Editorial note.} No omissions, hints or 'immediate' steps in the counted span.\par\medskip
\noindent\textit{Allowed standard prerequisites (editorial).}\begin{itemize}
\item Limit topology generators. Limits, generator lemmas of the chapter
\end{itemize}\medskip
\ifdefined\hypertarget\hypertarget{sourcepointer-1}{}\fi
\noindent\textit{Source reference.} Original proposition 01ZA: sources/stacks-chow/tags.txt, source lines 2585--2585. Cross-chapter standard prerequisite cited by the author inside the counted span; it is used as a pointer and its own proof is not part of this item.\par\medskip
\ifdefined\hypertarget\hypertarget{sourcepointer-2}{}\fi
\noindent\textit{Source reference.} Original reference 01TX: sources/stacks-chow/tags.txt, source lines 2390--2390. Cross-chapter standard prerequisite cited by the author inside the counted span; it is used as a pointer and its own proof is not part of this item.\par\medskip
\ifdefined\hypertarget\hypertarget{sourcepointer-3}{}\fi
\noindent\textit{Source reference.} Original reference 01TS: sources/stacks-chow/tags.txt, source lines 2385--2385. Cross-chapter standard prerequisite cited by the author inside the counted span; it is used as a pointer and its own proof is not part of this item.\par\medskip
\ifdefined\hypertarget\hypertarget{sourcepointer-4}{}\fi
\noindent\textit{Source reference.} Original reference 01T2: sources/stacks-chow/tags.txt, source lines 2359--2359. Cross-chapter standard prerequisite cited by the author inside the counted span; it is used as a pointer and its own proof is not part of this item.\par\medskip
\ifdefined\hypertarget\hypertarget{sourcepointer-5}{}\fi
\noindent\textit{Source reference.} Original lemma 01Z4: sources/stacks-chow/tags.txt, source lines 2579--2579. Cross-chapter standard prerequisite cited by the author inside the counted span; it is used as a pointer and its own proof is not part of this item.\par\medskip
\ifdefined\hypertarget\hypertarget{sourcepointer-6}{}\fi
\noindent\textit{Source reference.} Original lemma 01Z6: sources/stacks-chow/tags.txt, source lines 2581--2581. Cross-chapter standard prerequisite cited by the author inside the counted span; it is used as a pointer and its own proof is not part of this item.\par\medskip
\ifdefined\hypertarget\hypertarget{sourcepointer-7}{}\fi
\noindent\textit{Source reference.} Original lemma 01Z0: sources/stacks-chow/tags.txt, source lines 2575--2575. Cross-chapter standard prerequisite cited by the author inside the counted span; it is used as a pointer and its own proof is not part of this item.\par\medskip
\begin{lemma}
\label{lemma-locally-finite-type-in-finite-presentation}
Let $f : X \to S$ be a morphism of schemes.
Assume:
\begin{enumerate}
\item The morphism $f$ is locally of finite type.
\item The scheme $X$ is quasi-compact and quasi-separated.
\end{enumerate}
Then there exists a morphism of finite presentation
$f' : X' \to S$ and an immersion $X \to X'$ of schemes over $S$.
\end{lemma}

\begin{proof}
By
Proposition \ref{proposition-approximate}
we can write
$X = \lim_i X_i$ with each $X_i$ of finite type over $\mathbf{Z}$ and
with transition morphisms $f_{ii'} : X_i \to X_{i'}$ affine.
Consider the commutative diagram
$$
\xymatrix{
X \ar[r] \ar[rd] & X_{i, S} \ar[r] \ar[d] & X_i \ar[d] \\
& S \ar[r] & \Spec(\mathbf{Z})
}
$$
Note that $X_i$ is of finite presentation over $\Spec(\mathbf{Z})$, see
Morphisms,
Lemma \ref{morphisms-lemma-noetherian-finite-type-finite-presentation}.
Hence the base change $X_{i, S} \to S$ is of finite presentation by
Morphisms, Lemma \ref{morphisms-lemma-base-change-finite-presentation}.
Thus it suffices to show that the arrow $X \to X_{i, S}$ is an
immersion for $i$ sufficiently large.

\medskip\noindent
To do this we choose a finite affine open covering
$X = V_1 \cup \ldots \cup V_n$ such that
$f$ maps each $V_j$ into an affine open $U_j \subset S$.
Let $h_{j, a} \in \mathcal{O}_X(V_j)$ be a finite
set of elements which generate $\mathcal{O}_X(V_j)$ as
an $\mathcal{O}_S(U_j)$-algebra, see
Morphisms, Lemma \ref{morphisms-lemma-locally-finite-type-characterize}.
By Lemmas \ref{lemma-descend-opens} and \ref{lemma-limit-affine}
(after possibly shrinking $I$) we may assume that
there exist affine open coverings
$X_i = V_{1, i} \cup \ldots \cup V_{n, i}$
compatible with transition maps such that $V_j = \lim_i V_{j, i}$.
By Lemma \ref{lemma-descend-section} we can choose $i$ so large that each
$h_{j, a}$ comes from an element
$h_{j, a, i} \in \mathcal{O}_{X_i}(V_{j, i})$.
Thus the arrow in
$$
V_j \longrightarrow U_j \times_{\Spec(\mathbf{Z})} V_{j, i} =
(V_{j, i})_{U_j} \subset (V_{j, i})_S \subset X_{i, S}
$$
is a closed immersion. Since $\bigcup (V_{j, i})_{U_j}$ forms an open of
$X_{i, S}$ and since the inverse image of $(V_{j, i})_{U_j}$ in $X$
is $V_j$ it follows that $X \to X_{i, S}$ is an immersion.
\end{proof}

\end{document}
